\section{Benchmark 危机：从可比较分数到真实场景信号}

这场对谈最扎实的一段，是关于 benchmark 和 evaluation。谷雨的起手式很清楚：学术界和工业界需要的 evaluation 本来就不是同一种东西。学术界希望有可控实验来验证假设、定义新问题、比较方法能力；工业界则更关心评测上的提升能否迁移到真实部署，是否真的让用户完成任务、满意、留存或付费。

这一区分重要，是因为当下很多人一边说 benchmark 分数不重要，一边又离不开 benchmark。谢天宝给了一个更细的判断：benchmark 不是没用，而是要看处在什么阶段。在 0 到 1 阶段，它能定义问题、暴露短板、证明现有模型不会；一旦进入 60 到 99 的优化阶段，公开分数就越来越容易受数据合成、训练集过拟合和任务定制影响。这时只看分数，反而会误判真实可用性。

\begin{importantbox}{Benchmark 的位置变化}
Benchmark 最适合扮演 signal，而不是 final answer。它可以告诉社区“这里有一个问题”“模型在这里有短板”“某种能力开始出现”；但它不能单独回答“这个 Agent 在真实产品里是否好用”。越接近产业落地，越需要真实用户、真实 API、真实 case 和持续反馈闭环来校准。
\end{importantbox}

尚晏仪从旅行产品讲 reward，补上了工业界最难的一环。旅行产品最终可以看“有没有卖出一张票”，但这个 reward 离模型的每一步决策太远。用户没买票，可能是价格不合适，可能是路线不喜欢，可能是不信任 AI，也可能只是临时改变计划。用户说“不错”，也不等于真的会购买。也就是说，最终结果很清楚，但过程归因很模糊；主观反馈很多，但可训练信号很少。

这和学术 benchmark 的困难不同。学术任务里 reward 往往是预先定义的：做对得分，做错扣分。但真实产品里，用户自己也未必知道什么体验才是好体验。例如语音旅行助手什么时候该打断，什么时候该倾听，什么时候该长篇解释，这些问题不可能靠研究者凭直觉写一个 benchmark 就解决。尚晏仪说他们更依赖 A/B Test 和可持续试验基础设施，其实是在承认：用户偏好不是写出来的，是在真实交付中逐渐显影的。

\begin{knowledgebox}{0 到 1 与 60 到 99}
\begin{tabularx}{\textwidth}{>{\raggedright\arraybackslash}p{0.2\textwidth} Y Y}
\toprule
阶段 & Benchmark 的价值 & 主要风险 \\
\midrule
0 到 1 & 定义新任务，证明模型不会，给领域建立共同语言。 & 任务可能脱离真实场景，成为“为了 benchmark 而 benchmark”。 \\
60 到 99 & 作为局部能力信号，帮助工程团队定位短板。 & 分数容易被训练、数据增强、任务定制和选择性展示污染。 \\
真实产品 & 需要和用户反馈、商业结果、失败归因一起使用。 & reward 稀疏、主观、长链路，难以直接转成训练信号。 \\
\bottomrule
\end{tabularx}
\end{knowledgebox}

这一段背后的更大问题，是 Agent 研究开始从“能力是否存在”走向“能力如何稳定产生价值”。在前者里，benchmark 是一种科学仪器；在后者里，benchmark 只是产品运营系统的一部分。真正难的不是做出一个分数，而是把分数、用户反馈、人工分析、A/B Test、模型选择、prompt 策略、工具流程和业务指标连成循环。

\begin{warningbox}{不要把产品失败归因给一个单一模型分数}
Agent 产品失败时，很可能不是“模型不够强”这么简单。也许是任务定义错了，也许是 reward 太远，也许是交互不可信，也许是用户根本不想以这种方式完成任务。Benchmark 只能照亮一部分原因，不能替代产品归因。
\end{warningbox}

\subsection{本章小结}

Benchmark 危机的本质不是分数失去意义，而是分数的解释边界变窄了。学术界仍需要 benchmark 定义问题，工业界则必须把 evaluation 放进真实产品闭环里。GUI Agent 的评价尤其困难，因为它面对长链路任务、主观体验和稀疏 reward；真正有价值的 infrastructure，是能不断试、不断归因、不断把用户反馈转化为系统改进的 infrastructure。

